% ==============================================================================
% Course Syllabus: MSQF 531 Applied Time Series Analysis and Forecasting
% Formatted to fit on a single, clean official syllabus page
% ==============================================================================

\newpage
\thispagestyle{plain}

\begin{center}
    {\large\textbf{MSQF 531: APPLIED TIME SERIES ANALYSIS AND FORECASTING}} \hfill {\large\textbf{CREDITS: 4}}
\end{center}
\vspace{-0.6em}
\hrule height 0.8pt
\vspace{0.6em}

\noindent\textbf{Course Objectives:}
\begin{itemize}[leftmargin=1.5em, label=$\bullet$, noitemsep, topsep=1pt, parsep=1pt]
    \item \textit{Providing a clear explanation of the fundamental theory of time series analysis and forecasting.}
    \item \textit{The course features treatments of forecast improvement with regression and autoregression combination models and model and forecast evaluation, along with a sample size analysis for common time series models to attain adequate statistical power.}
\end{itemize}

\vspace{0.4em}
\noindent\textbf{Course Outcomes (COs):}
\begin{description}[leftmargin=1.4cm, labelsep=0.4em, itemsep=1pt, topsep=1pt, parsep=0pt, font=\normalfont\textbf]
    \item[CO1:] Understand and interpret the fundamental concepts and components of time series data
    \item[CO2:] Apply smoothing techniques for real world data using R and Excel
    \item[CO3:] Distinguish between stationary and non-stationary time series
    \item[CO4:] Application and modeling time-varying volatility using financial data
    \item[CO5:] Evaluate and compare various forecasting methods using real time data
\end{description}

\vspace{0.4em}
\noindent\textbf{Unit I: Introduction to Time Series}\\
\small
Definition and examples of Time Series Models --- Graphical Representation of Time Series Data --- Approaches of Time Series --- Additive and multiplicative approach --- Components and various decompositions of Time Series Models --- Numerical description of Time Series --- Data transformations --- Methods of estimation --- Trend, Seasonal and exponential.
\normalsize

\vspace{0.3em}
\noindent\textbf{Unit II: Smoothing Techniques and Univariate Time Series Models}\\
\small
Smoothing Techniques: Moving Averages: Simple, centered, double and weighted moving averages; single and double exponential smoothing --- Holt's and Winter's methods --- Exponential smoothing techniques for series with trend and seasonality --- First and Second order AR and MA Models --- ARMA/ARIMA Models --- Mixed ARMA/ARIMA models their statistical Properties --- Box-Jenkins methodology ACF and PACF functions --- Finite order $\AR(p)$ and $\MA(q)$ models.
\normalsize

\vspace{0.3em}
\noindent\textbf{Unit III: Multivariate and Multiple Equation Models}\\
\small
Trend stationary --- Stationary Time Series Models --- General Unit Root Tests: Dickey-Fuller Test, Augmented Dickey-Fuller Test --- Johansen Test for cointegration --- Granger causality --- Error Correction Model (ECM) --- Vector Autoregressive (VAR) model --- Vector Error Correction Model (VECM), estimation of Lag models.
\normalsize

\vspace{0.3em}
\noindent\textbf{Unit IV: Modeling Volatility}\\
\small
Impulse response function, variance decomposition --- Definition and representation of ARCH and GARCH Models --- Their use in financial time series data --- Volatility forecasting using $\GARCH(1,1)$ Model --- Diagnostic checking of model --- Analysis of residuals.
\normalsize

\vspace{0.3em}
\noindent\textbf{Unit V: Applications and Forecast Evaluation}\\
\small
Introduction to business forecasting --- Scope --- Types of forecasting --- Forecasting cycle --- Different forecasting techniques --- Exploring data patterns and choosing forecasting technique --- Managing forecasting process --- Measuring forecasting error --- Forecasting error comparison --- Forecast accuracy measures (MAE, MSE, RMSE, MAPE) --- Model validation and backtesting --- Scenario analysis and stress testing in financial forecasting --- Case studies: stock prices, interest rates, macroeconomic forecasting --- Use of software tools (R, Python) for applied time series modeling.
\normalsize

\vspace{0.5em}
\noindent\textbf{Books for Reference:}
\begin{enumerate}[leftmargin=1.5em, label=\arabic*., itemsep=1pt, topsep=1pt, parsep=0pt]
    \small
    \item Hyndman, R. J., \& Athanasopoulos, G. (2021). \textit{Forecasting: Principles and Practice} (3rd ed.).
    \item Hamilton, J. D. (2023). \textit{Time Series Analysis} (2nd ed.). Princeton University Press.
    \item Box, G. E. P., Jenkins, G. M., \& Reinsel, G. C. (1994). \textit{Time Series Analysis: Forecasting and Control} (3rd ed.). Prentice Hall.
    \item Makridakis, S., \& Wheelwright, S. C. (1997). \textit{Forecasting: Methods and Applications} (3rd ed.). John Wiley \& Sons.
    \item Mills, T. C. (1997). \textit{The Econometric Modelling of Financial Time Series} (2nd ed.). Cambridge University Press.
    \item Shumway, R. H., \& Stoffer, D. S. (2023). \textit{Time Series Analysis and Its Applications: With R Examples} (4th ed.). Springer.
    \item Tsay, R. S. (2024). \textit{Analysis of Financial Time Series} (4th ed.). John Wiley \& Sons.
\end{enumerate}

\vspace{0.4em}
\begin{center}
\footnotesize
\begin{tabular}{|c|c|c|c|c|c|}
\hline
\textbf{COs} & \textbf{PO1} & \textbf{PO2} & \textbf{PO3} & \textbf{PO4} & \textbf{PO5} \\
\hline
\textbf{CO1} & XXX & XX & XX & XX & XX \\
\hline
\textbf{CO2} & XXX & XXX & XX & XX & XX \\
\hline
\textbf{CO3} & XXX & XX & XX & XX & XX \\
\hline
\textbf{CO4} & XXX & XXX & XXX & XX & XX \\
\hline
\textbf{CO5} & XXX & XXX & XXX & XX & XXX \\
\hline
\end{tabular}
\end{center}
\normalsize

\newpage
